\section{The plan on one page}
\label{sec:summary}

\textbf{Goal.} A full-marks attempt at all 45 points and at every line of the brief's High
Distinction list. Every item in this plan is tied to a rubric line, a line of the brief, or a
promise made in Assignment 1 or 2. An item that earns none of these was left out.

\textbf{The marks, from the rubric.}

\begin{center}
\small
\begin{tabularx}{\textwidth}{@{}Y C{0.9cm} Y C{0.9cm}@{}}
\toprule
\textbf{Functionality} & \textbf{18} & \textbf{Visualisation} & \textbf{11} \\
\midrule
1. AI model integration & 3 & 1. Chart diversity and relevance (three or more types) & 3 \\
2. Front-end and back-end communication & 4 & 2. Interactivity (zoom, filtering, tooltips, updates) & 4 \\
3. Error handling and input validation & 3 & 3. Clarity (labels, legends, titles, units) & 2 \\
4. Back-end API design (REST, two or more methods) & 4 & 4. Performance and loading & 2 \\
5. Real-time updates and responsiveness & 4 & & \\
\midrule
\textbf{Source code} & \textbf{5} & \textbf{Usability} & \textbf{6} \\
\midrule
1. Execution (no errors or warnings, quick start) & 2 & 1. User interface (responsive, advanced components) & 3 \\
2. Comments & 1 & 2. User experience (feedback, personal touches, transitions) & 3 \\
3. Code structure & 2 & & \\
\midrule
\textbf{Report} (PDF, 8 pages at most) & \textbf{5} & & \\
\midrule
1. Formatting and Harvard referencing & 1 & 3. API endpoint documentation & 2 \\
2. System architecture and implementation & 2 & & \\
\bottomrule
\end{tabularx}
\end{center}

Forty of the 45 points are judged from the running application, and the brief says a video that
cannot be opened scores zero for most components. The application and its demonstration video
therefore come first; the report is five points.

\textbf{What will be built.}
\begin{enumerate}
  \item \textbf{A FastAPI back end around Assignment 2's \code{Predictor}.} It keeps the five
    endpoints Assignment 1 planned, with the same JSON, and adds a what-if prediction, a response
    curve, exports, a method list, a health check and the settings the front end needs
    (Section~\ref{sec:api}). Eleven endpoints are Must; eleven more come with the Should and Could
    items.
  \item \textbf{The React front end, grown from the Assignment 1 prototype.} Its three screens move
    from illustrative data to the real snapshot of 28 addresses and 7,966 transactions. Five views
    are added: a what-if lab, a models and limits page, a watchlist, an about page and a live-lookup
    dialog (Section~\ref{sec:frontend}).
  \item \textbf{User input to prediction in real time.} A validated form takes a hypothetical
    transaction (amount, direction, time) and shows every model's verdict as the user types. One
    request scores a whole range of amounts: measured at 0.16 s for seven amounts, identical to
    seven single calls that took 1.88 s (Section~\ref{sec:model}).
  \item \textbf{Interactive charts in Plotly.js.} Four are Must and cover three chart types (bar,
    scatter, line); two more are Should; three are Could (Section~\ref{sec:charts}). Each has zoom,
    tooltips, a legend, units, the snapshot date, a table view and an export.
  \item \textbf{Validation and error handling on both tiers}, from one table of rules and one
    catalogue of states (Section~\ref{sec:errors}).
  \item \textbf{The High Distinction features}: a method selector, a watchlist with notes, the
    guarded Explorer Assistant, a replay of an address's history, and a live lookup of any Bitcoin
    address, which needs change request CR-08.
  \item \textbf{The submission}: the report (8 pages), the code zip with its README, the video
    (7 minutes), the minutes for Weeks 1 to 12 and the contribution form (Section~\ref{sec:report}).
\end{enumerate}

\textbf{How this keeps the earlier promises.} Assignment 1's success measure was that AC-01 to
AC-03 and FR-05 pass on every snapshot address in the Assignment 3 build and are shown in the video.
That is milestone MS2. The wording rule does not change: the application says ``unusual'', ``flag''
and ``worth a look'', and every flagged item carries the line \emph{``Unusual activity requires
investigation; it does not indicate wrongdoing.''} Assignment 2's numbers are not touched: the
models are loaded as they were frozen at \code{results-v1.1}.

\textbf{The dates that matter.}
\begin{center}
\small
\begin{tabularx}{\textwidth}{@{}L{3.2cm}Y@{}}
\toprule
\textbf{Date} & \textbf{Milestone} \\
\midrule
Wed 11 Nov & CR-07 and CR-08 sent to the team; the questions in Appendix~\ref{app:cr} sent to the tutor. \\
Sun 15 Nov & MS1, the walking skeleton: the five Assignment 1 endpoints answer from the real snapshot and the three prototype screens show them. \\
Sun 22 Nov & MS2, the Must tier: the minimum viable demonstration passes on all 28 addresses, with the what-if lab, four charts, validation, error states and export. \\
Fri 27 Nov & MS3, feature freeze: the Should items that are finished stay; the rest are cut and recorded. \\
Sun 29 Nov & MS4, code freeze: a teammate installs from the README on Windows and on macOS; tag \code{v1.0-rc}. \\
Sat 5 Dec & MS5: everything submitted on Canvas. Sunday 6 December is spare. \\
\bottomrule
\end{tabularx}
\end{center}

\textbf{The critical path.} API contract → back-end endpoints → front end on the real API →
what-if and charts → MS2 → usability walkthrough and fixes → clean install → video → submission.
The Should items hang off MS2 and are never on this path. Assignment 1's baseline started this work
on about 28 October; it starts on 10 November, so 34 planned days are done in 26. Section
\ref{sec:schedule} shows how, and Section~\ref{sec:cutorder} says what is cut first.

\subsection{What needs the team this week}
\label{sec:needsteam}

\begin{team}
\begin{enumerate}
  \item \textbf{Generative AI (the tutor).} The brief says generative AI ``must not be used to
    generate content for this assignment'', allows it for learning, idea generation and debugging,
    and requires all use to be acknowledged in the report. For Assignment 2 the tutor allowed AI on
    condition of a separate declaration. Ask Arthur Hoang by Wed 11 Nov whether that arrangement
    covers Assignment 3, and build nothing that depends on the answer before it arrives
    (risk 32). The use log continues in \fp{docs/AI_USE_LOG.md} either way.
  \item \textbf{CR-07 and CR-08 (all four members).} Assignment 1's change control needs every
    member's recorded agreement to add the optional items and to allow a live lookup
    (Appendix~\ref{app:cr}).
  \item \textbf{Owners (weekly meeting, Week 9).} Section~\ref{sec:wp} proposes owners from
    Assignment 1's WBS. Each member needs an area they build and can explain in the video and to the
    tutor. Check the split against the contribution form's scaling rule before anyone starts.
  \item \textbf{Which commit was submitted for Assignment 2 (Cao Tri Tran).} The repository has no
    \code{v1.0} tag. Tag the submitted commit so the back end can name the exact models it loads.
  \item \textbf{The Week 9--12 slides (anyone).} Put them in the shared folder when they are
    released; decisions D-A3-2 to D-A3-4 are checked against them.
\end{enumerate}
\end{team}
